\documentclass[11pt]{article}
\usepackage[margin=1in]{geometry}
\usepackage{amsmath,amssymb,amsthm}
\usepackage{array,longtable,booktabs}
\usepackage[hidelinks]{hyperref}

\newtheorem{theorem}{Theorem}
\newtheorem{proposition}[theorem]{Proposition}
\newtheorem{corollary}[theorem]{Corollary}
\newtheorem{example}[theorem]{Example}
\newcommand{\Ann}{\operatorname{Ann}}
\newcommand{\vol}{\operatorname{vol}}

\title{P6-T04: Source Density and Tick Carriers Do Not Yet Select Operational Metric Scale}
\author{The AEther-Flow Research Project}
\date{2026-07-26}

\begin{document}
\maketitle

\begin{center}
\fbox{\parbox{0.92\textwidth}{
\textbf{Status and authority.}
This is a task-local \emph{draft/control} result about proposal-only
source-extension carriers.  It proves a conditional volume-gauge theorem and
a precise candidate-scoped nonselection obstruction.  It is not a
canonical-ontology edit, source-law adoption or rejection, physical clock or
rod construction, effective-metric derivation, Gate Chair verdict, global
no-go theorem, publication, or completed derivation.
}}
\end{center}

\section{Question and fixed input}

P6-T03 leaves as its complete source-side geometric input
\[
  \mathcal D_1(S,V)=\bigl([V]_+,\Ann(V),\preceq_V\bigr),
\]
where \(S\) is a smooth four-dimensional source arena and \(V\) is the
proposal-only oriented nonzero vector field from P6-T02.  It proves that this
one-ray datum does not select a Lorentzian conformal class.

P6-T04 asks whether source-defined measure or operational carriers can
nevertheless fix scale and measured durations and lengths.  We test the
stronger package
\[
 \mathcal D_{\rm cal}
 =\bigl(\mathcal D_1,\mu,\{\theta_A\}_{A\in\mathcal A}\bigr),
\]
where \(\mu\) is a smooth positive source density and each \(\theta_A\) is a
dimensionless phase carrier.  These are deliberately favorable
\emph{proposal-only test carriers}; current ontology does not derive their
physical meaning, normalization, compatibility, or universality.  No target
atlas, target metric, GR null cone, proper time, detector response, matter
datum, or benchmark fit is admitted.

\section{Assumption and dimensional audit}

\begin{longtable}{>{\raggedright\arraybackslash}p{0.10\textwidth}
                  >{\raggedright\arraybackslash}p{0.38\textwidth}
                  >{\raggedright\arraybackslash}p{0.41\textwidth}}
\toprule
ID & Test assumption & Status and dimensional consequence\\
\midrule
\endhead
A1 & \(S\) is smooth and four dimensional, and \(V\) is smooth, nonzero, and
positively oriented. & Inherited candidate scope; not a physical spacetime or
clock conclusion.\\
A2 & Only \(\mathcal D_1\), a positive density \(\mu\), and scalar phases
\(\theta_A\) are supplied. & Exact enlarged input under test; no hidden
metric, coframe, connection, or target calibration.\\
A3 & \(\mu\) is treated as a source density rather than a metric volume form. &
Its physical dimension and relation to spacetime volume are not derived.\\
A4 & Each \(\theta_A\) is dimensionless and a tick is a \(2\pi\) phase
increment. & Tick counts are meaningful conditionally; no seconds-per-tick
constant is supplied.\\
A5 & Positive reparameterizations \(V\mapsto fV\) preserve the oriented ray and
unparameterized forward flow. & An absolute flow parameter is not part of
\(\mathcal D_1\).\\
A6 & The historical scoped source-extension \(g_{\rm eff}\) object is not a
premise and its status is unchanged. & No protected scope expansion or
unscoped physical metric is inferred.\\
A7 & No conversion constant with dimensions of time, length, or volume is
available. & A numerical density or phase count cannot alone establish a
physical unit.\\
A8 & No target-side geometry, matter response, detector behavior, or exact-GR
benchmark data are premises. & Target fitting and benchmark calibration are
excluded.\\
\bottomrule
\end{longtable}

\section{What a density can fix conditionally}

\begin{theorem}[Conditional volume gauge]
\label{thm:volume-gauge}
Let \([g]\) be an independently supplied Lorentzian conformal class on a
four-dimensional oriented region and let \(\mu\) be a positive density.
There is a unique representative \(g_\mu\in[g]\) whose absolute metric volume
density equals \(\mu\).
\end{theorem}

\begin{proof}
Choose any representative \(g\in[g]\).  Under \(g\mapsto g'=\Omega^2g\)
with \(\Omega>0\), the four-dimensional volume density transforms as
\[
  |\vol_{g'}|=\Omega^4|\vol_g|.
\]
Thus
\[
  \Omega=\left(\frac{\mu}{|\vol_g|}\right)^{1/4}
\]
is positive and uniquely enforces \(|\vol_{\Omega^2g}|=\mu\).  Replacing the
initial representative by \(\widehat g=\omega^2g\) replaces the displayed
factor by \(\widehat\Omega=\Omega/\omega\), so the resulting metric is
independent of the initial representative.
\end{proof}

\begin{corollary}
A density is a conditional Weyl-gauge selector, not a constructor of the
missing conformal shape.  The theorem cannot be invoked on
\(\mathcal D_1\) alone because P6-T03 proves that \(\mathcal D_1\) supplies no
unique \([g]\).  Physical use would additionally require a source law
identifying \(\mu\) with operational spacetime volume and fixing its units.
\end{corollary}

\section{Fixed-density anisotropy countermodel}

\begin{theorem}[Unimodular density, cone, and rod nonselection]
\label{thm:unimodular}
In local flow-box coordinates \((s,y^1,y^2,y^3)\) with \(V=\partial_s\), the
one-parameter family
\[
  g_a=-ds^2+e^{2a}(dy^1)^2+e^{-2a}(dy^2)^2+(dy^3)^2,
  \qquad a\in\mathbb R,
\]
has the same oriented \(V\)-ray, annihilator, forward flow preorder, and
absolute volume density.  Distinct values of \(a\) are nevertheless
nonconformal, generally have different null cones, and assign different
transverse rod lengths.
\end{theorem}

\begin{proof}
Every \(g_a\) is Lorentzian and has \(g_a(V,V)=-1\).  The ray,
\(\Ann(V)\), and flow preorder depend only on \(V\).  Moreover,
\[
 \det(g_a)=-e^{2a}e^{-2a}=-1,
\]
so every member has the same density
\(|ds\wedge dy^1\wedge dy^2\wedge dy^3|\).
If \(g_b=\Omega^2g_a\), evaluation on \(V\) gives \(\Omega^2=1\); evaluation
on \(\partial_{y^1}\) then gives \(e^{2b}=e^{2a}\), hence \(b=a\).
For \(a=0\), \(V+\partial_{y^1}\) is null, while for \(a\ne0\) its squared
norm is \(-1+e^{2a}\).  A coordinate segment of \(y^1\)-length \(L\) has
comparison length \(e^aL\), while a \(y^2\)-segment has \(e^{-a}L\).
Thus fixed density controls the product of transverse scales but not their
shape or rod comparisons.
\end{proof}

\section{Clock phases and device agreement}

For a phase carrier \(\theta_A\), define the local phase rate
\(\omega_A=V(\theta_A)\).  Under the positive reparameterization
\(V_f=fV\),
\[
  \omega_A\longmapsto \omega_A^{(f)}=f\,\omega_A.
\]
Therefore neither \(\omega_A\) nor a duration inferred as
\(2\pi/\omega_A\) is invariant without a normalization of \(V\).
For two nonzero phase carriers, however,
\[
   R_{AB}=\frac{V(\theta_A)}{V(\theta_B)}
\]
is invariant under \(V\mapsto fV\).  This supports a bounded
\emph{relative-device agreement protocol}: two devices agree on a connected
test region when \(R_{AB}\) is a source-derived constant and controlled
perturbations keep it within a declared error band.  The protocol compares
cycle ratios only.  It does not supply an absolute physical time unit, and
current ontology derives neither the constant ratio nor its robustness law.

The countermodel of Theorem~\ref{thm:unimodular} may carry the same phases
\(\theta_A\) for every \(a\).  Hence equal tick data and equal density coexist
with different null cones and rod lengths.  A phase normalization can at most
constrain a selected direction conditionally; it does not reconstruct the
missing transverse geometry.

\section{Precise obstruction}

\begin{proposition}[Volume-tick calibration nonselection]
\label{prop:obstruction}
There is no single-valued operational Lorentzian metric reconstruction from
only \(\mathcal D_{\rm cal}\) as defined above.  Identical one-ray, density,
and phase-carrier data admit the pair \(g_0,g_a\) with \(a\ne0\), which have
different conformal classes, null cones, and rod lengths.  Positive
reparameterization of \(V\) separately leaves the oriented ray and
unparameterized flow unchanged while rescaling every absolute phase rate.
\end{proposition}

The exact disposition is
\texttt{OBST-P6T04-VOLUME-TICK-CALIBRATION-NONSELECTION-001}.  The current
ontology does not derive:
\begin{enumerate}
\item a Lorentzian conformal shape to which the conditional volume-gauge
theorem could be applied;
\item a source law identifying a positive density as physical volume with a
dimensionful normalization;
\item an absolute clock-frequency or flow-parameter normalization;
\item a transverse rod/coframe comparison law; or
\item a compatibility and robustness theorem forcing independent clock and
rod devices to agree.
\end{enumerate}
Supplying any of these as a conversion rule, coframe, constitutive response,
or normalized measure would be a new source-side primitive unless separately
derived.

This is not a theorem that the project can never derive operational metric
scale.  Same-milestone continuation remains open for materially richer
source-side input, and protected human authority remains mandatory before
any canonical-ontology candidate, adoption, rejection, physical metric,
clock, rod, or benchmark claim.  Current adoption is
\texttt{blocked\_adoption\_open\_continuation}.

\section{Distance-to-GR disposition}

P6-T04 is mathematically decisive within the enlarged input tested here.  It
proves one conditional positive theorem and an explicit fixed-density,
fixed-tick countermodel.  The exact unchanged route from one-ray plus
unnormalized density and phase carriers to an operational Lorentzian metric
is locally frozen.  The historical scoped \(g_{\rm eff}\) object remains
unchanged.  Physical conformal geometry, metric scale, clocks, rods,
\(M_{\rm src}\), matter coupling, Einstein equations, benchmark promotion,
proof authority, publication, push, and completed derivation remain blocked.

\section{Project source note}

The fixed project basis is the registered P6-T03 conformal-nonselection
theorem, the canonical target-side geometry dictionary, the historical scoped
source-extension \(g_{\rm eff}\) candidate and Gate Chair record, the v21
implementation plan, the Distance-to-GR burden map, and handoff-0880.  These
sources are held at the exact hashes recorded in the accompanying
machine-readable specification.  The scaling, determinant,
reparameterization, and countermodel calculations above are reproduced
directly; no external physical datum is used.

\end{document}
